\documentclass[12pt]{article}
\usepackage[utf8]{inputenc}
\usepackage[T1]{fontenc}
\usepackage[letterpaper, margin=0.75in]{geometry}
\usepackage{hyperref}
\usepackage{enumitem}
\usepackage{graphicx}

% Set paragraph formatting for a block letter style
\setlength{\parindent}{0pt}
\setlength{\parskip}{1em}

\begin{document}
% --- Logo Placement ---
% This places the image at the top center. 
% Ensure your file is named 'logo.png' in the Overleaf sidebar.
\begin{flushright}
    \includegraphics[width=0.5\textwidth]{logo.png}
\end{flushright}

% ---------- Sender ----------
{\large\textbf{Rubén Rellán-Álvarez}}\\
Associate Professor\\
Department of Molecular and Structural Biochemistry\\
N.C. Plant Sciences Initiative\\
North Carolina State University, Raleigh, NC, USA\\
\href{mailto:rrellan@ncsu.edu}{rrellan@ncsu.edu}

% --- Date ---
\today

% --- Recipient Block ---
Editor-in-Chief\\
\textit{The Plant Genome}\\
Crop Science Society of America
 
\vspace{1.5em}
Dear Dr. Henry Nguyen,
 
\vspace{0.5em}
We are pleased to submit our manuscript, ``The Sorghum Lipid Database (SoLD): population-scale lipidomics linking environmental and genetic variation in the Sorghum Association Panel,'' for consideration for publication in \textit{The Plant Genome}.
 
Population-scale lipidomics resources remain scarce in crops relative to genomic and transcriptomic datasets, even though membrane and storage lipids are central to nutrient acclimation and abiotic-stress tolerance. We present SoLD, a curated lipidomics resource for the Sorghum Association Panel (SAP)---a widely used community diversity panel---profiled by high-resolution LC--MS across two contrasting field regimes: a nutrient-sufficient control and a reduced-input regime with lower nitrogen and phosphorus and earlier planting. After QC-based normalization (SERRF) and spatial correction (SpATS), we quantified 244 lipid species and resolved reproducible, environment-associated remodeling of the sorghum lipidome, including sulfoquinovosyldiacylglycerol depletion, triacylglycerol enrichment, phosphatidylserine-centered phospholipid redistribution, and coordinated lysophospholipid changes.
 
We integrated these phenotypes with genome-wide association analyses of individual lipid species, lipid-class sums, and class ratios in each environment, using recurrence across traits to prioritize candidate loci. These loci connect lipid variation to nutrient-status signaling, cell-wall remodeling, developmental regulation, and cold-associated barrier formation. Lipid-reaction network (LINEX2) and lipid-ontology (LION) analyses further place candidate genes within a biochemical framework, defining a focused set of hypotheses for future functional validation.
 
We believe this work is well aligned with \textit{The Plant Genome}'s focus on genetic and genomic resources for crop species. Beyond the dataset and analyses, the study delivers an openly accessible, interactive database---the SoLD Shiny application---together with a public code and data repository, allowing the community to move from lipid classes to individual molecular species and their associated candidate loci for hypothesis generation, comparative analyses, and target prioritization. Because the two cohorts were grown in different field years, we interpret all condition contrasts conservatively as between-trial differences and present the GWAS results as prioritized candidates rather than validated causal genes, positioning SoLD as a community resource and integrative framework.
 
This manuscript is original, has not been published previously, and is not under consideration for publication elsewhere; a preprint is available on bioRxiv (\href{https://doi.org/10.64898/2026.06.05.727974}{doi:10.64898/2026.06.05.727974}). All authors have read and approved the manuscript and agree to its submission. All data and code underlying the study are publicly available, with links provided in the manuscript.
 
Thank you for considering our manuscript. We look forward to your evaluation.
 
\vspace{1em}
Sincerely,

% --- Signature Block ---
\textbf{Nirwan Tandukar, Rub\'en Rell\'an-\'Alvarez} \\
Department of Molecular and Structural Biochemistry \\
N.C. Plant Sciences Initiative \\
North Carolina State University, Raleigh, NC, USA \\
\href{mailto:ntanduk@ncsu.edu}{ntanduk@ncsu.edu}, \href{mailto:rrellan@ncsu.edu}{rrellan@ncsu.edu}

\end{document}